% Options for packages loaded elsewhere
\PassOptionsToPackage{unicode}{hyperref}
\PassOptionsToPackage{hyphens}{url}
%
\documentclass[
  10pt,
  letterpaper,
]{article}
\usepackage{amsmath,amssymb}
\usepackage{iftex}
\ifPDFTeX
  \usepackage[T1]{fontenc}
  \usepackage[utf8]{inputenc}
  \usepackage{textcomp} % provide euro and other symbols
\else % if luatex or xetex
  \usepackage{unicode-math} % this also loads fontspec
  \defaultfontfeatures{Scale=MatchLowercase}
  \defaultfontfeatures[\rmfamily]{Ligatures=TeX,Scale=1}
\fi
\usepackage{lmodern}
\ifPDFTeX\else
  % xetex/luatex font selection
\fi
% Use upquote if available, for straight quotes in verbatim environments
\IfFileExists{upquote.sty}{\usepackage{upquote}}{}
\IfFileExists{microtype.sty}{% use microtype if available
  \usepackage[]{microtype}
  \UseMicrotypeSet[protrusion]{basicmath} % disable protrusion for tt fonts
}{}
\makeatletter
\@ifundefined{KOMAClassName}{% if non-KOMA class
  \IfFileExists{parskip.sty}{%
    \usepackage{parskip}
  }{% else
    \setlength{\parindent}{0pt}
    \setlength{\parskip}{6pt plus 2pt minus 1pt}}
}{% if KOMA class
  \KOMAoptions{parskip=half}}
\makeatother
\usepackage{xcolor}
\usepackage[margin=1in]{geometry}
\usepackage{color}
\usepackage{fancyvrb}
\newcommand{\VerbBar}{|}
\newcommand{\VERB}{\Verb[commandchars=\\\{\}]}
\DefineVerbatimEnvironment{Highlighting}{Verbatim}{commandchars=\\\{\}}
% Add ',fontsize=\small' for more characters per line
\newenvironment{Shaded}{}{}
\newcommand{\AlertTok}[1]{\textcolor[rgb]{1.00,0.00,0.00}{\textbf{#1}}}
\newcommand{\AnnotationTok}[1]{\textcolor[rgb]{0.38,0.63,0.69}{\textbf{\textit{#1}}}}
\newcommand{\AttributeTok}[1]{\textcolor[rgb]{0.49,0.56,0.16}{#1}}
\newcommand{\BaseNTok}[1]{\textcolor[rgb]{0.25,0.63,0.44}{#1}}
\newcommand{\BuiltInTok}[1]{\textcolor[rgb]{0.00,0.50,0.00}{#1}}
\newcommand{\CharTok}[1]{\textcolor[rgb]{0.25,0.44,0.63}{#1}}
\newcommand{\CommentTok}[1]{\textcolor[rgb]{0.38,0.63,0.69}{\textit{#1}}}
\newcommand{\CommentVarTok}[1]{\textcolor[rgb]{0.38,0.63,0.69}{\textbf{\textit{#1}}}}
\newcommand{\ConstantTok}[1]{\textcolor[rgb]{0.53,0.00,0.00}{#1}}
\newcommand{\ControlFlowTok}[1]{\textcolor[rgb]{0.00,0.44,0.13}{\textbf{#1}}}
\newcommand{\DataTypeTok}[1]{\textcolor[rgb]{0.56,0.13,0.00}{#1}}
\newcommand{\DecValTok}[1]{\textcolor[rgb]{0.25,0.63,0.44}{#1}}
\newcommand{\DocumentationTok}[1]{\textcolor[rgb]{0.73,0.13,0.13}{\textit{#1}}}
\newcommand{\ErrorTok}[1]{\textcolor[rgb]{1.00,0.00,0.00}{\textbf{#1}}}
\newcommand{\ExtensionTok}[1]{#1}
\newcommand{\FloatTok}[1]{\textcolor[rgb]{0.25,0.63,0.44}{#1}}
\newcommand{\FunctionTok}[1]{\textcolor[rgb]{0.02,0.16,0.49}{#1}}
\newcommand{\ImportTok}[1]{\textcolor[rgb]{0.00,0.50,0.00}{\textbf{#1}}}
\newcommand{\InformationTok}[1]{\textcolor[rgb]{0.38,0.63,0.69}{\textbf{\textit{#1}}}}
\newcommand{\KeywordTok}[1]{\textcolor[rgb]{0.00,0.44,0.13}{\textbf{#1}}}
\newcommand{\NormalTok}[1]{#1}
\newcommand{\OperatorTok}[1]{\textcolor[rgb]{0.40,0.40,0.40}{#1}}
\newcommand{\OtherTok}[1]{\textcolor[rgb]{0.00,0.44,0.13}{#1}}
\newcommand{\PreprocessorTok}[1]{\textcolor[rgb]{0.74,0.48,0.00}{#1}}
\newcommand{\RegionMarkerTok}[1]{#1}
\newcommand{\SpecialCharTok}[1]{\textcolor[rgb]{0.25,0.44,0.63}{#1}}
\newcommand{\SpecialStringTok}[1]{\textcolor[rgb]{0.73,0.40,0.53}{#1}}
\newcommand{\StringTok}[1]{\textcolor[rgb]{0.25,0.44,0.63}{#1}}
\newcommand{\VariableTok}[1]{\textcolor[rgb]{0.10,0.09,0.49}{#1}}
\newcommand{\VerbatimStringTok}[1]{\textcolor[rgb]{0.25,0.44,0.63}{#1}}
\newcommand{\WarningTok}[1]{\textcolor[rgb]{0.38,0.63,0.69}{\textbf{\textit{#1}}}}
\setlength{\emergencystretch}{3em} % prevent overfull lines
\providecommand{\tightlist}{%
  \setlength{\itemsep}{0pt}\setlength{\parskip}{0pt}}
\setcounter{secnumdepth}{-\maxdimen} % remove section numbering
\usepackage{xurl}
\setlength{\emergencystretch}{3em}
\usepackage{microtype}
\ifLuaTeX
  \usepackage{selnolig}  % disable illegal ligatures
\fi
\usepackage{bookmark}
\IfFileExists{xurl.sty}{\usepackage{xurl}}{} % add URL line breaks if available
\urlstyle{same}
\hypersetup{
  hidelinks,
  pdfcreator={LaTeX via pandoc}}

\author{}
\date{}

\begin{document}

\section{DDTA research work plan after documentation-layer
closure}\label{ddta-research-work-plan-after-documentation-layer-closure}

\textbf{WORK PLAN - REVISION 9}

\textbf{Status:} active BA2 execution plan after BA2-T1
proposition-structure lower-bound derivation.

\textbf{Supersedes:} Revision 8 only for forward execution state; R1-R8
remain historical research records.

\subsection{1. Fixed prior state}\label{fixed-prior-state}

\begin{itemize}
\tightlist
\item
  Chapters 2-4: CLOSED / FINAL for current thesis scope.
\item
  Documentation layer: CLOSED.
\item
  BA0-R systems-modeling prior-art research: CLOSED.
\item
  BA0 responsibility/non-goals: CLOSED.
\item
  BA1 minimal BAE identity ontology: CLOSED.
\item
  \texttt{BAReferent} and \texttt{BAProposition}: ACCEPTED first-class
  semantic identity families.
\item
  ThreatForge is a case study/tooling instrument, not DDTA semantic
  authority.
\end{itemize}

\subsection{2. Closed BA1 identity
boundary}\label{closed-ba1-identity-boundary}

BA1 accepts exactly:

\begin{Shaded}
\begin{Highlighting}[]
\NormalTok{BAReferent}
\NormalTok{BAProposition}
\end{Highlighting}
\end{Shaded}

No \texttt{Action}, \texttt{Relation}, \texttt{Participation},
\texttt{Role}, \texttt{Behavior}, \texttt{Information}, \texttt{Store},
\texttt{Contract}, \texttt{Boundary} or \texttt{State} identity family
may be added by BA2 unless a concrete counterexample satisfies the BA1
reopen criteria.

\subsection{3. BA2 objective}\label{ba2-objective}

BA2 defines the smallest stable methodology-neutral
proposition/relation/action semantics over the BA1 identities while
preserving:

\begin{Shaded}
\begin{Highlighting}[]
\NormalTok{semantic identity family}
\NormalTok{    !=}
\NormalTok{semantic operator / participation role / modifier}
\NormalTok{    !=}
\NormalTok{lexical label / authoring synonym}
\end{Highlighting}
\end{Shaded}

BA2 must support mechanical selection/projection without importing
threat-method categories or requiring consumers to reinterpret raw
prose.

\subsection{4. BA2-T1 - proposition shape and participation-role lower
bound}\label{ba2-t1---proposition-shape-and-participation-role-lower-bound}

\textbf{Status: COMPLETED / PROVISIONAL CANDIDATE.}

BA2-T1 tested universal binary SPO, positional n-ary tuples and fixed
domain/method slots against facial-access and order-fulfillment
evidence.

\subsubsection{Result}\label{result}

\begin{Shaded}
\begin{Highlighting}[]
\NormalTok{Pure binary SPO as universal core       REJECTED}
\NormalTok{Binary propositions as special cases   SUPPORTED}
\NormalTok{Anonymous positional participation      REJECTED}
\NormalTok{Role{-}bound n{-}ary participation          FORCED CANDIDATE}
\NormalTok{Explicit semantic operator              FORCED CANDIDATE}
\NormalTok{Scoped semantic modifier capability     REQUIRED / ENCODING OPEN}
\NormalTok{Third BA1 identity family               NOT FORCED}
\end{Highlighting}
\end{Shaded}

Active lower-bound candidate:

\begin{Shaded}
\begin{Highlighting}[]
\NormalTok{BAProposition}
\NormalTok{|{-} semanticOperator                  1}
\NormalTok{|{-} participation                     1..*}
\NormalTok{|    |{-} role                         1}
\NormalTok{|    \textasciigrave{}{-} term                         1}
\NormalTok{\textasciigrave{}{-} scopedModifier                    0..*}
\end{Highlighting}
\end{Shaded}

At least one participant term is a BAReferent. Typed local values may
appear when independent semantic identity is not required.

Referent classification remains separate from contextual proposition
role. BA2-T1 supports explicit method-neutral classification capability
but does not freeze its vocabulary or cardinality.

\subsection{5. BA2-T2 - operator, role and modifier vocabulary pressure
test}\label{ba2-t2---operator-role-and-modifier-vocabulary-pressure-test}

\textbf{NEXT / NOT EXECUTED BY THIS PLAN REVISION.}

BA2-T2 must attempt to derive the smallest reusable methodology-neutral
vocabulary needed to instantiate the T1 structure.

Pressure areas:

\begin{enumerate}
\def\labelenumi{\arabic{enumi}.}
\tightlist
\item
  operator families needed by both corpora without copying document
  predicate labels mechanically;
\item
  role semantics for actor/executor, source/destination,
  information/object, request/result, correlation/context,
  provider/consumer/responsibility and affected state;
\item
  whether roles are globally canonical or operator-scoped;
\item
  operator-role compatibility/cardinality constraints;
\item
  representation boundary between role-bound participants and scoped
  modifiers;
\item
  condition, state, failure, applicability, realization,
  atomicity/concurrency and negation/polarity pressure;
\item
  reusable referent classification needed for mechanical human/method
  projections;
\item
  bounded method-consumer mapping without STRIDE/DFD leakage.
\end{enumerate}

\subsubsection{Exit condition}\label{exit-condition}

Produce a falsifiable vocabulary candidate and explicit
rejected/deferred alternatives. Do not close BA2 if important
operator/role/modifier ambiguity remains.

\subsection{6. Later BA2 closure sequence - provisional
only}\label{later-ba2-closure-sequence---provisional-only}

If BA2-T2 yields a stable vocabulary candidate, later microsteps should
pressure-test:

\begin{itemize}
\tightlist
\item
  cross-corpus regression;
\item
  lexical canonical label versus synonym separation;
\item
  role/cardinality consistency;
\item
  whether any relation/action semantic identity needs a registry-level
  stable key;
\item
  BA2 closure and reopen criteria.
\end{itemize}

The exact number of later microsteps is evidence-driven, not
precommitted.

\subsection{7. BA3 - derivation, provenance and
authority}\label{ba3---derivation-provenance-and-authority}

Only after BA2 closure, BA3 must close:

\begin{itemize}
\tightlist
\item
  document-to-BAReferent/BAProposition derivation rules;
\item
  source locators;
\item
  grounded, derived and diagnostic/unresolved materialization;
\item
  identity/equivalence/lifecycle rules;
\item
  acceptance/rejection/stale semantics;
\item
  source-to-analysis and analysis-to-source traceability.
\end{itemize}

Governed documentation remains project authority.

\subsection{8. BA4 - projections}\label{ba4---projections}

Define derived human and method projections from the same accepted BA
semantics. No view becomes source authority.

\subsection{9. BA5 - vocabulary and optional
assistance}\label{ba5---vocabulary-and-optional-assistance}

Define lexical controlled vocabulary, authoring synonyms and optional
assistance only after semantic operator/role contracts are closed.
NLP/LLM proposals remain candidate-producing support and cannot
establish project semantics silently.

\subsection{10. BA6 - Base Analysis closure and
regression}\label{ba6---base-analysis-closure-and-regression}

Regress the implemented Base Analysis design against closed corpora and
at least one structurally different holdout. BA6 closes the complete
Base Analysis design.

\subsection{11. Analysis envelope remains
downstream}\label{analysis-envelope-remains-downstream}

Only after Base Analysis closure:

\begin{itemize}
\tightlist
\item
  A1 - AnalysisRecord / execution envelope;
\item
  A2 - common reviewed Finding boundary;
\item
  A3 - change/provenance integration;
\item
  formal methodology overlays and STRIDE/STRIDE-AI evaluations.
\end{itemize}

The earlier STRIDE transfer probe remains only a bounded consumer probe.

\subsection{12. Next authorized
microstep}\label{next-authorized-microstep}

After the BA2-T1 package is reviewed, committed and remotely verified,
execute only:

\begin{quote}
\textbf{BA2-T2 - semantic operator, participation-role and
scoped-modifier vocabulary pressure test.}
\end{quote}

Do not start BA3, formal STRIDE overlay design, Common Finding schema or
implementation work.

\end{document}
